
The aim of this work was to develop and evaluate a short-term PM2.5 forecasting system for Almaty built entirely on open data and deployable on free-tier infrastructure. The five objectives stated in the Introduction have been addressed in full.

\textbf{Objective 1 - Data ingest and curation.} A reproducible SQLite-backed pipeline was implemented across \texttt{src/ingest/} and \texttt{src/db.py} to fetch hourly PM2.5 observations from 192 OpenAQ LCS stations and merge them with seven Open-Meteo meteorological variables and CAMS-Global PM2.5 hindcasts across the 562-day training window (2024-10-01 \ensuremath{\rightarrow} 2026-04-15). The raw archive contains 578,993 PM2.5 sensor-hour records from 190 sensors. After plausibility filtering (0--500 \textmu{}g/m$^3$) and a minimum-sensor-count requirement of three active stations per hour, the archive reduces to 11,553 city-wide hourly medians. Meteorological data (13,488 continuous hourly rows) and CAMS data (13,488 rows) join to this series without loss. The pipeline is idempotent: it skips existing Parquet snapshots and can be rebuilt from scratch by any user with a valid OpenAQ API key, satisfying the reproducibility requirement central to the open-science framing of this work.

\textbf{Objective 2 - Feature engineering.} Twenty-seven features were constructed across six families: five autoregressive PM2.5 lags (1, 3, 6, 12, 24 h), three rolling means (3, 6, 24 h, anchored at t -- 1 to prevent leakage), six meteorological variables including raw and log-transformed boundary-layer height and its hour-over-hour delta, wind Cartesian decomposition, temperature gradient, six cyclical calendar encodings, and a binary heating-season indicator for October 15 through March 31. The feature set was designed by domain knowledge rather than automated selection, producing an interpretable 341:1 sample-to-feature ratio that ensures the subsequent model comparison is not confounded by dimensionality effects.

\textbf{Objective 3 - Model training and comparison.} XGBoost, Ridge regression, persistence, and CAMS were evaluated on a 2,311-hour held-out test set coinciding with the winter 2025--26 peak season. XGBoost achieves the best performance at the 6-hour horizon (RMSE = 23.44 \textmu{}g/m$^3$, R$^2$ = 0.545), beating persistence by 22.2\% and CAMS by 41.1\% in RMSE. At the 12 and 24-hour horizons, Ridge regression outperforms XGBoost (RMSE 25.85 vs.~28.67 \textmu{}g/m$^3$ at +12 h; 27.24 vs.~31.48 \textmu{}g/m$^3$ at +24 h) - a result attributed to the degradation of non-linear feature interactions with forecast lead time and to the relatively limited size of the training set. CAMS registers negative R$^2$ at all three horizons, confirming that a globally coarse chemistry-transport model cannot resolve Almaty's sub-grid inversion dynamics. Every locally trained model beats CAMS by at least 16\% in RMSE across all horizons and model types. This consistent margin supports the primary scientific claim: locally trained, observation-driven forecasting adds reproducible value over the best freely available global alternative for this geography.

\textbf{Objective 4 - Error regime analysis.} Two distinct failure modes were characterised. During inversion-onset events, all lag-based models lag the observed PM2.5 ramp by approximately 2--3 hours; at a typical ramp rate of 10--20 \textmu{}g/m$^3$/h, this corresponds to a peak-concentration underestimate of 20--60 \textmu{}g/m$^3$. The XGBoost model partially compensates through the BLH delta feature, which provides an advance meteorological signal that pure autoregressive models lack. During rapid clearing events, autoregressive models over-predict PM2.5 for 2--4 hours after the boundary-layer breakdown. The BLH threshold analysis in \S{}3.4.5 shows that the model's RMSE in the low-BLH regime (BLH \textless{} 100 m) is 3.3 times higher than in the well-mixed regime (BLH \textgreater{} 500 m), quantifying the concentration of forecast difficulty in the physically demanding inversion-onset scenario.

\textbf{Objective 5 - Streamlit web application.} A single-screen Streamlit application (\texttt{src/app/main.py}) was implemented with a Folium sensor map colour-coded by the AQI scale, a forecast chart showing the +6/+12/+24 h XGBoost predictions with \ensuremath{\pm}1 MAE confidence bands, and a district-level summary table with trend indicators. The application refreshes predictions by fetching the latest 48 hours from the OpenAQ and Open-Meteo APIs and running inference on the three pre-trained models (each approximately 2 MB) in approximately 30--60 seconds on a standard network connection.

\textbf{Scientific novelty.} Both novelty claims stated in the Introduction are supported by the experimental results. The present work delivers the first publicly described and evaluated PM2.5 forecasting model specifically calibrated for Almaty's LCS network and inversion-dominated winter meteorology; no prior peer-reviewed study has benchmarked a locally trained gradient-boosted model against CAMS for this geography. The finding that Ridge regression outperforms XGBoost at 12 and 24-hour horizons in an LCS-based, moderate-size urban dataset also contributes to the methodological debate on model complexity selection for data-scarce Central Asian urban environments, where reference-grade training targets and long historical records are not available.

\textbf{Broader implications.} The results establish general principles applicable beyond the Almaty case. Locally trained forecasting outperforms global CTM output by at least 16\% in RMSE and this gap holds across all horizons, motivating investment in data infrastructure and model development in other Central Asian cities. Model complexity should be matched to horizon and dataset size: XGBoost dominates at short horizons where non-linear interactions provide signal, while Ridge regression is more reliable at longer horizons where the training sample is insufficient to support complex non-linear fits. The open, parameterised pipeline developed in this work provides a direct template for researchers or agencies in Shymkent, Karaganda, Bishkek, or Tashkent to replicate the forecasting system for their local networks within hours, requiring only a change of coordinates and API window in \texttt{src/config.py}.

\textbf{Reliability assessment and limitations.} The reliability of these conclusions is bounded by three acknowledged limitations. The PM2.5 training targets are uncalibrated LCS readings subject to humidity-induced positive biases of 20--60\% at relative humidity above 80\%, which are conditions frequently occurring during Almaty's winter inversions. The reported RMSE values are relative to this biased target, not to a reference-grade mass concentration standard. The training window of 562 days spans approximately 1.3 heating seasons; published guidance on gradient-boosted PM2.5 models suggests that 2--3 full annual cycles improve inter-annual generalisation, meaning the model has not been exposed to the full range of synoptic variability that Almaty experiences across consecutive winters. No formal hyperparameter optimisation was applied to XGBoost; the reported reversal of model ranking between XGBoost and Ridge regression at longer horizons may shift if a Bayesian search over the hyperparameter space were conducted, though the fundamental conclusion that Ridge regression is competitive at 12 and 24-hour horizons is expected to hold given the dataset size.

\textbf{Contribution to open science in Central Asia.} The complete pipeline - API ingest, SQLite storage, feature construction, model training, evaluation, and Streamlit deployment - is version-controlled in a single Python repository with a fully locked dependency specification (\texttt{uv.lock}) and \texttt{pyproject.toml}. Any researcher can rebuild the dataset, retrain the models, and redeploy the application using only a personal OpenAQ API key and a laptop with internet access. The pipeline's configuration is centralised in \texttt{src/config.py}, where the city centre coordinates, training window, and API parameters are set; changing these four values is the only modification required to adapt the pipeline to another Central Asian city. This design choice - prioritising portability over technical sophistication - reflects a deliberate positioning of the work as a practical baseline for the region rather than a methodological benchmark for the global forecasting literature.

Future work should prioritise two directions. The first is LCS bias correction via a short co-location campaign alongside a Kazhydromet reference instrument, which would improve the absolute accuracy of the training targets and enable comparison with published RMSE figures from reference-grade studies. The second is integration of the Open-Meteo forecast BLH product - forecast rather than analysed boundary-layer height at the prediction time - which would extend the model's look-ahead on inversion evolution and partially compensate for the 2--3 hour lag at onset transitions that is the principal operational weakness of the current architecture. A third direction of practical value is a weighted ensemble combining XGBoost's 6-hour precision with Ridge regression's 12 and 24-hour stability, which is straightforward to implement and expected to match or exceed both individual models at all horizons without additional training data. These three improvements would move the system from its current research-prototype status toward an operationally deployable public health tool for one of Central Asia's most polluted cities.
